\documentclass{beamer}

\usetheme{Madrid}
\usecolortheme{default}

\title{Counterfactual Validation and Exploitation of Longitudinal Treatment Data}
\author{Malek Ben Salah}
\institute{Institut Gustave Roussy / CentraleSupélec}
\date{\today}

\usepackage{amsmath}
\usepackage{graphicx}
\usepackage{tikz}
\usepackage{booktabs}

% Placeholder figure command
\newcommand{\placeholderfig}[2]{
\begin{figure}
\centering
\fbox{
\begin{minipage}[c][0.25\textheight][c]{0.8\textwidth}
\centering
\textbf{#1} \\[0.5em]
\textit{#2}
\end{minipage}}
\end{figure}
}

\begin{document}

\begin{frame}
\titlepage
\end{frame}

%------------------------------------------------
\begin{frame}{Motivation and Experimental Goals}
\begin{itemize}
    \item Electronic health records aggregated at treatment-line level
    \item Objective: estimate \textbf{counterfactual survival} under alternative treatment sequences
    \item Key challenge: validating counterfactual predictions without ground-truth outcomes
\end{itemize}

\vspace{0.5em}
\textbf{Experimental objectives:}
\begin{enumerate}
    \item Validate counterfactual survival predictions
    \item Assess robustness to confounding and temporal structure
    \item Prepare foundations for a treatment recommendation policy
\end{enumerate}
\end{frame}

%------------------------------------------------
\begin{frame}{Modeling Setup (Reminder)}
\begin{itemize}
    \item Sequential model of treatment lines using LSTM
    \item Patient representation updated at each treatment start
    \item Treatment-specific MLP heads predicting survival distributions
    \item Adversarial loss to enforce balance w.r.t. propensity scores
\end{itemize}

\placeholderfig
{Model architecture overview}
{LSTM over treatment lines with treatment-specific survival heads}
\end{frame}

%------------------------------------------------
\begin{frame}{Experimental Roadmap}
\begin{enumerate}
    \item Factual consistency and internal counterfactual checks
    \item Representation balance and deconfounding diagnostics
    \item Counterfactual ordering and temporal coherence
    \item Off-policy evaluation of induced recommendation rules
    \item Robustness and sensitivity analyses
\end{enumerate}
\end{frame}

%------------------------------------------------
\begin{frame}{Factual Consistency Check}
\textbf{Question:}  
Does the model recover observed outcomes when intervening on the factual treatment?

\begin{itemize}
    \item Compare predicted vs observed survival under factual treatment
    \item Stratify by treatment line and propensity score quintiles
\end{itemize}

\placeholderfig
{Calibration of factual survival predictions}
{Observed vs predicted survival curves, stratified by treatment}
\end{frame}

%------------------------------------------------
\begin{frame}{Counterfactual Placebo Test}
\textbf{Question:}  
Does the model hallucinate treatment effects under randomized interventions?

\begin{itemize}
    \item Shuffle treatment assignments within propensity strata
    \item Re-train counterfactual model
    \item Expect collapse of estimated treatment effects
\end{itemize}

\placeholderfig
{Placebo counterfactual effects}
{Distribution of estimated treatment effects after label shuffling}
\end{frame}

%------------------------------------------------
\begin{frame}{Representation Balance Diagnostics}
\textbf{Question:}  
Is treatment assignment predictable from learned patient representations?

\begin{itemize}
    \item Train post-hoc classifier to predict treatment from latent state
    \item Compare with and without adversarial deconfounding
\end{itemize}

\placeholderfig
{Treatment predictability from latent representations}
{AUROC of treatment classifier across treatment lines}
\end{frame}

%------------------------------------------------
\begin{frame}{Pairwise Counterfactual Stability}
\textbf{Question:}  
Are pairwise treatment comparisons stable for similar patients?

\begin{itemize}
    \item Compute pairwise survival differences
    \[
    \Delta_{a,b}(x_t) = \hat{S}(a|x_t) - \hat{S}(b|x_t)
    \]
    \item Evaluate variance within propensity-matched groups
\end{itemize}

\placeholderfig
{Pairwise treatment effect stability}
{Variance of pairwise effects across matched patient strata}
\end{frame}

%------------------------------------------------
\begin{frame}{Clinical Plausibility Checks}
\textbf{Question:}  
Do estimated effects align with known clinical patterns?

\begin{itemize}
    \item Stratification by biomarkers (HR, HER2, etc.)
    \item Check monotonicity and sign consistency
\end{itemize}

\placeholderfig
{Counterfactual survival by biomarker subgroup}
{Predicted survival curves under different treatments}
\end{frame}

%------------------------------------------------
\begin{frame}{Temporal Coherence Across Lines}
\textbf{Question:}  
Are counterfactual effects coherent across successive treatment lines?

\begin{itemize}
    \item Compare predicted treatment effects at line $t$ vs $t+1$
    \item Expect smooth degradation, not abrupt reversals
\end{itemize}

\placeholderfig
{Temporal evolution of counterfactual effects}
{Treatment effect magnitude as a function of line number}
\end{frame}

%------------------------------------------------
\begin{frame}{History Ablation Study}
\textbf{Question:}  
How sensitive are predictions to treatment history?

\begin{itemize}
    \item Mask parts of prior treatment history
    \item Measure deviation in counterfactual survival estimates
\end{itemize}

\placeholderfig
{History sensitivity analysis}
{Change in predicted survival vs length of ablated history}
\end{frame}

%------------------------------------------------
\begin{frame}{From Counterfactuals to Policy Evaluation}
\textbf{Policy definition:}
\[
\pi(x_t) = \arg\max_a \hat{S}(a|x_t)
\]

\begin{itemize}
    \item Evaluate induced policy using off-policy estimators
    \item Compare with observed clinical policy
\end{itemize}

\placeholderfig
{Estimated policy value}
{Overall survival under learned vs observed policy}
\end{frame}

%------------------------------------------------
\begin{frame}{Retrospective Regret Analysis}
\textbf{Question:}  
Where does the model disagree with observed clinical choices?

\begin{itemize}
    \item Compute regret per treatment line
    \[
    r_t = \max_a \hat{S}(a|x_t) - \hat{S}(a_{\text{obs}}|x_t)
    \]
    \item Stratify by risk and line number
\end{itemize}

\placeholderfig
{Regret distribution}
{Estimated regret across treatment lines}
\end{frame}

%------------------------------------------------
\begin{frame}{Sensitivity to Unmeasured Confounding}
\textbf{Question:}  
How strong must hidden confounding be to change conclusions?

\begin{itemize}
    \item Inject synthetic confounders
    \item Measure tipping points for treatment ranking changes
\end{itemize}

\placeholderfig
{Sensitivity analysis}
{Treatment ranking stability vs confounder strength}
\end{frame}

%------------------------------------------------
\begin{frame}{Summary and Next Steps}
\begin{itemize}
    \item Extensive counterfactual validation without ground truth
    \item Evidence for stability, plausibility, and robustness
    \item Natural extension toward dynamic treatment recommendation
\end{itemize}

\vspace{0.5em}
\textbf{Next steps:}
\begin{itemize}
    \item Look-ahead planning over treatment sequences
    \item Constraint-aware recommendation rules
    \item Prospective validation design
\end{itemize}
\end{frame}

%------------------------------------------------

\begin{frame}{Longitudinal Survival Modeling: The Challenge}
\begin{itemize}
    \item Goal: Predict survival at multiple time points during patient follow-up
    \item Available data:
    \begin{itemize}
        \item Intermediate time points with evolving features
        \item Final outcome (event or censored)
    \end{itemize}
    \item \textbf{Key Question}: How should we label intermediate time points?
\end{itemize}

\vspace{0.3cm}
\begin{block}{Two Competing Approaches}
\begin{enumerate}
    \item \textbf{Ultimate Outcome}: Label all time points with final outcome
    \item \textbf{Censored}: Label only with event status at that specific time
\end{enumerate}
\end{block}
\end{frame}

\begin{frame}{Approach Comparison}
\begin{columns}[T]
\column{0.48\textwidth}
\begin{block}{Ultimate Outcome Approach}
\small
\textbf{Advantages:}
\begin{itemize}
    \item Captures early warning signals
    \item Predicts "who will die eventually"
    \item Stable predictions over time
    \item Simpler interpretation
\end{itemize}

\vspace{0.2cm}
\textbf{Risks:}
\begin{itemize}
    \item Information leakage from future
    \item Ignores time-varying risk
    \item Overly optimistic performance
\end{itemize}
\end{block}

\column{0.48\textwidth}
\begin{block}{Censored Approach}
\small
\textbf{Advantages:}
\begin{itemize}
    \item Proper survival modeling
    \item No future information leakage
    \item Learns time-varying risk
    \item Realistic performance estimates
\end{itemize}

\vspace{0.2cm}
\textbf{Limitations:}
\begin{itemize}
    \item May miss subtle early signals
    \item More complex interpretation
    \item Requires more data
\end{itemize}
\end{block}
\end{columns}
\end{frame}

\begin{frame}{Data Structure Example}
\textbf{Patient who dies at time step 5:}

\begin{table}[h]
\centering
\small
\begin{tabular}{|c|c|c|c|}
\hline
\textbf{Time} & \textbf{Features} & \textbf{Ultimate} & \textbf{Censored} \\
\hline
1 & $X_1$ & Death (1) & Alive (0) \\
2 & $X_2$ & Death (1) & Alive (0) \\
3 & $X_3$ & Death (1) & Alive (0) \\
4 & $X_4$ & Death (1) & Alive (0) \\
5 & $X_5$ & Death (1) & Death (1) \\
\hline
\end{tabular}
\end{table}

\vspace{0.3cm}
\begin{alertblock}{Key Insight}
The censored approach reflects \emph{actual information available} at each time point, while the ultimate outcome approach uses future knowledge.
\end{alertblock}
\end{frame}

\begin{frame}{Proposed Solution: Multi-Task Learning}
\textbf{Key Idea}: Combine both approaches using a residual-inspired architecture

\vspace{0.3cm}
\begin{itemize}
    \item \textbf{Main pathway}: Learn immediate time-specific risk (censored)
    \item \textbf{Skip connection}: Learn ultimate outcome trajectory
    \item \textbf{Shared encoder}: Both tasks learn from same representations
\end{itemize}

\vspace{0.3cm}
\begin{block}{Mathematical Formulation}
\begin{equation*}
\mathcal{L}_{\text{total}} = \alpha \cdot \mathcal{L}_{\text{immediate}} + \beta \cdot \mathcal{L}_{\text{ultimate}}
\end{equation*}
where:
\begin{itemize}
    \item $\mathcal{L}_{\text{immediate}}$: Binary cross-entropy for time-specific events
    \item $\mathcal{L}_{\text{ultimate}}$: Binary cross-entropy for final outcome
    \item $\alpha, \beta$: Task weights (e.g., $\alpha = 0.7, \beta = 0.3$)
\end{itemize}
\end{block}
\end{frame}

% \begin{frame}{Multi-Task Architecture}
% \begin{center}
% \includegraphics[width=0.95\textwidth]{multi_task_survival_diagram.pdf}
% \end{center}

% \vspace{0.2cm}
% \footnotesize
% \textbf{Note}: Branch 1 (green) is primary for clinical decisions; Branch 2 (red) serves as auxiliary task to improve feature learning
% \end{frame}

\begin{frame}{Benefits of Multi-Task Approach}
\begin{enumerate}
    \item \textbf{Richer Feature Learning}
    \begin{itemize}
        \item Model learns representations that satisfy both objectives
        \item Early biomarkers weighted appropriately for long-term prognosis
    \end{itemize}
    
    \item \textbf{Regularization Effect}
    \begin{itemize}
        \item Ultimate outcome task prevents overfitting to short-term noise
        \item Improved generalization
    \end{itemize}
    
    \item \textbf{Complementary Information}
    \begin{itemize}
        \item Immediate risk: "What's the danger now?"
        \item Ultimate outcome: "What's the overall trajectory?"
    \end{itemize}
    
    \item \textbf{Clinical Flexibility}
    \begin{itemize}
        \item Use Branch 1 for time-specific decisions
        \item Use Branch 2 for long-term prognosis discussions
    \end{itemize}
\end{enumerate}
\end{frame}

% \begin{frame}{Implementation Considerations}
% \begin{block}{Hyperparameter Tuning}
% \begin{itemize}
%     \item Task weights: Start with $\alpha > \beta$ (e.g., 0.7:0.3)
%     \item Monitor both losses during training
%     \item Adjust if one task dominates
% \end{itemize}
% \end{block}

% \begin{block}{Evaluation Strategy}
% \begin{itemize}
%     \item \textbf{Branch 1}: Time-dependent C-index, Brier score at each time point
%     \item \textbf{Branch 2}: Overall AUC for ultimate outcome prediction
%     \item \textbf{Both}: Calibration curves at different time horizons
% \end{itemize}
% \end{block}

% \begin{alertblock}{Key Consideration}
% Ensure sufficient dataset size to support the added model complexity without overfitting
% \end{alertblock}
% \end{frame}

% \begin{frame}{Expected Outcomes}
% \begin{columns}[T]
% \column{0.48\textwidth}
% \textbf{What the Model Learns:}
% \begin{itemize}
%     \item Acute risk factors (infections, vital sign changes)
%     \item Chronic progression markers (organ failure)
%     \item Subtle early indicators of poor prognosis
% \end{itemize}

% \column{0.48\textwidth}
% \textbf{Clinical Applications:}
% \begin{itemize}
%     \item Real-time risk monitoring
%     \item Early intervention triggers
%     \item Prognostic counseling
%     \item Resource allocation
% \end{itemize}
% \end{columns}

% \vspace{0.5cm}
% \begin{block}{Hypothesis}
% The multi-task approach will outperform either single-task approach by:
% \begin{itemize}
%     \item Better capturing time-varying dynamics (vs. ultimate-only)
%     \item Identifying early risk signals (vs. censored-only)
% \end{itemize}
% \end{block}
% \end{frame}

% \begin{frame}{Next Steps}
% \begin{enumerate}
%     \item \textbf{Implementation}
%     \begin{itemize}
%         \item Build shared encoder (LSTM/Transformer for temporal features)
%         \item Implement dual output heads
%         \item Define combined loss function
%     \end{itemize}
    
%     \item \textbf{Validation}
%     \begin{itemize}
%         \item Compare against single-task baselines
%         \item Perform ablation studies on task weights
%         \item Test on held-out temporal validation set
%     \end{itemize}
    
%     \item \textbf{Clinical Validation}
%     \begin{itemize}
%         \item Assess calibration at multiple time horizons
%         \item Evaluate clinical utility (decision curve analysis)
%         \item Get clinician feedback on predictions
%     \end{itemize}
% \end{enumerate}
% \end{frame}


\end{document}